Los resultados obtenidos muestran diferencias numéricas observables entre las funciones de pérdida evaluadas, junto con un rendimiento limitado del sistema respecto al \textit{baseline}. Para interpretar correctamente los valores proporcionados, es importante distinguir entre las métricas y el tipo de información que aporta cada una:

\begin{itemize}
\item \textbf{SI-SDR}: mide la calidad de la separación obtenida.

\item \textbf{SI-SDRi}: indica la mejora respecto a una referencia base.

\item \textbf{SI-SAR}: permite observar la presencia de artefactos introducidos durante el proceso de separación, como ruido, distorsiones o presencia de señales de otras fuentes.

\end{itemize}

En el escenario de separación de dos fuentes, ninguna función supera el \textit{baseline}. Los mejores resultados corresponden a:

\begin{itemize}
\item \texttt{Deep-Feature-EMD}: con un valor de SI-SDRi medio de \(-0.081\).

\item \texttt{Deep-Feature}: con un valor de SI-SDRi medio de \(-0.091\).

\end{itemize}

Estos dos modelos son los más cercanos al \textit{baseline}, aunque sin llegar a superarlo. Sus valores de SI-SDRi están muy próximos a cero y, además, presentan los mejores valores de SI-SAR, con \(23.435\) y \(22.996\), respectivamente. Esto sugiere que se generan separaciones con menor presencia relativa de artefactos.

Por debajo de las funciones experimentales aparecen las pérdidas basadas en la norma \(L2\), especialmente \texttt{L2-Freq} y \texttt{L2}. En este escenario, \texttt{L2-Freq} obtiene un SI-SDRi medio de \(-2.262\), ligeramente mejor que \texttt{L2}, que alcanza \(-2.515\). Esto indica que la variante frecuencial mejora algo el resultado respecto a la versión básica en el escenario de dos fuentes.

Las funciones basadas en \(L1\), las variantes logarítmicas y \texttt{Mask-l1} obtienen los peores resultados. En particular, \texttt{logl1} y \texttt{logl2} presentan valores de SI-SDRi muy negativos, lo que indica un empeoramiento claro frente al \textit{baseline}. Esto sugiere que las pérdidas logarítmicas no han proporcionado una señal de optimización adecuada para mejorar la separación.

En el escenario de separación de cuatro fuentes, los resultados son diferentes. Una función supera el \textit{baseline}: \texttt{Deep-Feature-EMD}, con un SI-SDRi medio de \(0.006\). Esta mejora es positiva, pero extremadamente reducida, por lo que debe interpretarse como una mejora marginal y no como una diferencia sólida. El valor obtenido por esta función está muy próximo a cero y, por tanto, no permite afirmar que el modelo supere de forma amplia al sistema de referencia.

Las funciones que ocupan las mejores posiciones en el ranking de SI-SDR medio son:

\begin{itemize}
\item \texttt{Deep-Feature-EMD}: con un SI-SDR medio de \(-5.375\).

\item \texttt{Deep-Feature}: con un SI-SDR medio de \(-5.399\).

\item \texttt{L2}: con un SI-SDR medio de \(-5.445\).

\item \texttt{L2-Freq}: con un SI-SDR medio de \(-5.449\).

\end{itemize}

Resulta relevante observar el comportamiento del SI-SAR en cuatro fuentes. Las funciones \texttt{L2} y \texttt{L2-Freq} obtienen valores superiores a los de \texttt{Deep-Feature-EMD}, con \(20.263\) y \(20.437\), frente a \(17.562\). Esto puede interpretarse como que, aunque \texttt{Deep-Feature-EMD} obtiene el mejor resultado en términos de mejora de SI-SDR, las pérdidas basadas en \(L2\) generan separaciones con menor presencia relativa de artefactos.

Las pérdidas logarítmicas muestran un mal comportamiento en ambos escenarios. En cuatro fuentes, \texttt{logl1}, \texttt{Log-mag} y \texttt{Log-compressed-l2} se sitúan en las últimas posiciones del ranking y presentan valores negativos de SI-SAR, lo que indica una peor calidad de reconstrucción y una mayor introducción de artefactos. Este resultado sugiere que, para la arquitectura y el protocolo utilizados en este trabajo, estas pérdidas no se ajustan bien al objetivo de separación.

En conjunto, los resultados indican que las funciones experimentales basadas en características profundas son las más competitivas dentro del conjunto de modelos entrenados. En particular, \texttt{Deep-Feature-EMD} obtiene el mejor comportamiento global, al ser el único modelo que consigue superar, aunque de forma marginal, el \textit{baseline} en cuatro fuentes y al quedar muy próximo a él en el escenario de dos fuentes.

La principal conclusión del estudio es que el proyecto ha permitido implementar y validar un pipeline completo de entrenamiento y evaluación para la separación musical, pero los resultados cuantitativos muestran que los modelos entrenados todavía presentan limitaciones. Las mejoras respecto al \textit{baseline} son escasas, especialmente en el escenario de dos fuentes, y la única mejora positiva en cuatro fuentes es muy ligera.

Esto apunta a la necesidad de continuar ajustando la arquitectura, ampliar el entrenamiento, revisar las funciones de pérdida y estudiar estrategias adicionales para mejorar la calidad de separación.
